\begin{abstract}
% (0.5/0.6 vs. 0.4/0.46 accuracies across 2 models)
% We also provide an in-depth analysis of alternative complexity estimation techniques based on expert assessment via model-as-judge (MASJ) and chain-of-thought entropy aggregation with ROC AUC scores of 0.57 and 0.7 accordingly. 

General-purpose Large Language Models (LLMs) are frequently fine-tuned through supervised fine-tuning (SFT) to enhance performance in specific domains. Better results can be achieved by distilling the chain-of-thought of a larger model at the cost of numerous expensive calls and a much greater amount of data.
We propose a novel blueprint for efficient fine-tuning that uses reasoning only for complex data identified by entropy. Specifically, across three small open models ($\approx 3B$) we split the training data into complexity categories by a single token answer entropy (ROC AUC $0.73$), fine-tune large language models (LLMs) via SFT and distillation, and show that our pipeline significantly outperforms the standard SFT approach ($0.58$ vs $0.45$ average accuracy) and outperforms the distillation approach ($0.58$ vs $0.56$ average accuracy) while using $81\%$ less data.
We publish our code\footnote{\url{https://github.com/LabARSS/complexity-aware-fine-tuning}} and data\footnote{\url{https://github.com/LabARSS/complexity-aware-fine-tuning?tab=readme-ov-file\#data}} to facilitate further research in this direction.
\end{abstract}

% General purpose Large Language Models (LLMs) are frequently fine-tuned to improve performance in niche domains. Although fine-tuning is a standard practice, we still lack a deep understanding of how to aggregate data for better results. In this work, we show that the entropy-based output estimation provides a meaningful guideline for fine-tuning data preparation. Specifically, across two small open models ($\approx 3B$) we find that a single token answer entropy shows ROC AUC score of $\approx 0.73$ and allows us to split the training data into three complexity categories to apply different tuning mechanisms. As result, we propose a novel blueprint for efficient fine-tuning that outperforms the standard approach $0.55$ vs $0.43$ average accuracy. Our findings show immediate enhancements in fine-tuning performance. We publish our code\footnote{\url{https://github.com/LabARSS/complexity-aware-fine-tuning}} and data\footnote{\url{https://huggingface.co/datasets/LabARSS}} to facilitate further investigation and immersion of the numerical complexity analysis.
